\scanpage{40}

\uline{Proposition} (i) Si $(LG)$ est satisfaite, alors pour
\deleted{$\forall$ voisinage ouvert}
$\operatorname{Pic}(X)\longrightarrow\operatorname{Pic}(X_{/Y})$
est injectif, et plus généralement, on a un hom injectif
\[
(\alpha)\qquad
\mathop{\mathrm{lim}}_{\substack{\longrightarrow\\U\supset Y}}
\operatorname{Pic}(U)\longrightarrow\operatorname{Pic}(X_{/Y})
\]
avec hom de transition injectifs.

(ii) Si $(LG)$ strict est satisfait, alors $\alpha$ est un isom.

\uline{Corollaire} Supposons que $Y$ ``est'' un diviseur de
Cartier sur $X$, soit $\deleted{\unsure{\mathcal O}}\Lambda\simeq
\underline L(Y)$ le faisceau inv. associé, donc
$i^*(\Lambda)^{-1}\simeq J/J^{\unsure{2}}$,
où $J$ est l'Idéal définissant $Y$.
[$i:Y\longrightarrow X$ désigne l'injection canonique.

\margin{introduire aussi les $Y_n$ pour exprimer ces propriétés}

(i) Supposons $H^1(Y,\Lambda^{-n})=0$ pour $n=1,2,\ldots$,
et $(LG)$. Alors $\operatorname{Pic}(X)\mathbin{-}\operatorname{Pic}(Y)$
est \uline{injectif}.
[\unsure{plus gén.} \deleted{\unsure{et on a les mêmes}}
\deleted{$(\alpha)$ est \unsure{injectif}}
les $\operatorname{Pic}(U)\longrightarrow\operatorname{Pic}(Y)$
sont injectifs, \unsure{donc} \deleted{\unsure{surjectifs}}
\unsure{de même}

(ii) Supposons
$\unsure{\deleted{H^1(Y,\Lambda^{-n})}|H^1(Y,\Lambda^{-n})}=0$
\uuline{et} $H^2(Y,\Lambda^{-n})=0$ pour $n=1,2,\ldots$,
et $(LG)$ strict.
Alors
\[
\mathop{\mathrm{lim}}_{\substack{\longrightarrow\\U\supset Y}}
\operatorname{Pic}(U)\simeq\operatorname{Pic}(Y).
\]

\margin{\[
\mathop{\mathrm{lim}}_{\substack{\longrightarrow\\U\supset Y}}
\operatorname{Pic}(U)\longrightarrow\deleted{\operatorname{Pic}(Y)}
\]
[et hom de transition injectifs]}

\margin{39}
